\subsubsection{構造設計記述}

\term{構造設計記述}{Structural Design Description} は、ソフトウェアの静的な構造を表す。つまり、どの構成要素があり、それらがどう関係し、どの責務を持つかを示す。

代表的な記法は次のとおりである。

\par\smallskip\noindent\refstepcounter{table}\label{tab:ch03-04}\textbf{表 \thetable: 第03章 ソフトウェア設計 表04}\par\smallskip
\begin{longtable}{L{42mm}L{112mm}}
\toprule
\textbf{記法} & \textbf{説明} \\
\midrule
\endhead
クラス図・オブジェクト図 & クラス、オブジェクト、属性、操作、関連を表す。オブジェクト指向設計でよく使われる。 \\
コンポーネント図 & 交換可能な構成要素と、その提供・要求インタフェースを表す。 \\
CRCカード & クラス名、責務、協調相手をカード形式で整理する。初期設計や会話に向く。 \\
配置図 & ソフトウェアをどの物理ノードや実行環境へ配置するかを表す。 \\
実体関連図 & データの実体、属性、関係を表す。データ設計でよく使われる。 \\
インタフェース記述言語 & 構成要素が公開する操作名、型、引数、戻り値などを記述する。 \\
構造チャート & プログラムの呼び出し構造を表す。どのモジュールがどのモジュールを呼ぶかを示す。 \\
\bottomrule
\end{longtable}
